\section{Method}
\label{sec:method}

The method has two components.  \textbf{M1} maps the \emph{form} of a governing
PDE to a valid GP kernel for each tensor mode; \textbf{M2} fits a functional
Tucker model whose factors carry those kernels.  Nothing else is claimed: the
inference machinery is standard and is described only to make the comparison
reproducible.

\subsection{M1: from a PDE form to per-mode kernels}

\paragraph{Operator to joint spectrum.}
For a stationary operator driven by noise, $\Lop u = w$, the solution is a
stationary Gaussian process with spectral density
\begin{equation}
  \Sop(\wvec;\theta) = \bigl|\Lhat_\theta(\wvec)\bigr|^{-2}\, \Sw(\wvec).
  \label{eq:joint-spectrum}
\end{equation}
Only the functional form of $\Lhat$ is assumed known; $\theta$ is set to nominal
values, and \Cref{sec:experiments} measures how little that choice matters.

\paragraph{Nonnegative separation.}
$\Sop$ is not separable across modes, whereas a functional tensor needs a
one-dimensional function per mode.  We therefore project it,
\begin{equation}
  \min_{\lambda_q \ge 0,\; s_{dq} \ge 0}
  \Bigl\| \Sop - \sum_{q} \lambda_q\, s_{1q} \circ \cdots \circ s_{Dq} \Bigr\|_F^2 ,
  \label{eq:nncp}
\end{equation}
and report the relative separation error as an \emph{a priori} applicability
diagnostic, computed before any data is seen.

\paragraph{The eigenbasis depends on the boundary condition.}
A stationary operator's eigenbasis is fixed by its boundary condition, not by
its symbol alone.  On a periodic domain it is the complex exponential pair; on a
no-flux domain the Laplacian eigenfunctions are cosines,
$\psi_k(x) = \sqrt{2}\cos(\pi k x)$ with eigenvalues $(\pi k)^2$.  This matters
in practice rather than in principle: initial-value data is not periodic in
time, and forcing $f(0) = f(1)$ is a large unphysical constraint.
\todo{cite the 1.30 -> 1.08 measurement}

\paragraph{Validity.}
For any nonnegative spectrum $s$ and either basis,
\begin{equation}
  k_s(x,x') = \sum_k s(k)\, \psi_k(x)\, \psi_k(x')
  \label{eq:psd}
\end{equation}
is positive semi-definite, being a nonnegatively weighted sum of outer products;
so is any nonnegative combination of such spectra.  Every step from the
nonnegative separation to the final kernel therefore preserves validity.  In the
cosine basis the pointwise variance is not constant, which is correct: on a
bounded domain the covariance genuinely is non-stationary near the edges.

\subsection{M2: functional Tucker host}
\begin{equation}
  u(x_1,x_2,x_3) = \sum_{r_1,r_2,r_3} G_{r_1 r_2 r_3}\,
  f_{1 r_1}(x_1) f_{2 r_2}(x_2) f_{3 r_3}(x_3),
  \qquad
  f_{d r}(x) = \phi_{s_d}(x)^\top a_{dr},
  \label{eq:tucker}
\end{equation}
with per-mode ranks chosen independently.  A CP host is not sufficient: real
two-dimensional fields are fields of blobs, which are not $x \otimes y$
separable at any rank sparse observations could identify, even though their
multilinear rank is small.  Per-mode ranks matter concretely -- a field of
multilinear rank $[2,11,11]$ needs a $242$-entry core, while an equal rank of
$11$ would need $1331$.

Coefficients carry a mean-field Gaussian posterior, $q(a_{dr}) = \Norm(\mu_{dr},
\diag(\sigma_{dr}^2))$ against $p(a_{dr}) = \Norm(0,I)$, and the core, noise and
any routing weights are point estimates.  Training maximises the usual ELBO.
The coefficient budget is $2\sum_d R_d F$, independent of how many spectra a
bank contains, so comparisons between banks of different sizes are not
confounded by parameter count.
